\subsection{Positive self-intersection and a quantitative divisibility obstruction}
\label{subsec:positive-square-picard}

We prove the quantitative bound of Theorem~\ref{thm:positive-excess-picard}
using only positive
self-intersection of the distinguished curve. The ambient surface may
be normal and singular away from that curve. The argument keeps track
of the arithmetic genus of its inverse image under a normalized
$p$-root cover and of the loss of self-intersection needed to make that
inverse image smooth. This permits a finite sequence of root transfers
when the genus excess is sufficiently small.

\begin{lemma}\label{lem:positive-square-logarithmic}
Let $S$ be a normal integral projective surface over an algebraically
closed field $k$ of characteristic $p>0$. Let $C\subset S$ be a smooth
integral effective Cartier divisor contained in the smooth locus, with
$C^2>0$. If $B\ne\OO_S$ has order $p$ and $B|_C\simeq\OO_C$, its
logarithmic form $\theta$ is nonzero and has exact order $m\geq1$ along
$C$. Moreover $p\mid m$, and its first tangential coefficient is a
nonzero section of
\[
 \Omega_C^1\otimes\OO_C(-mC).
\]
In particular, $2g(C)-2\geq mC^2$. At equality that section is nowhere
zero.
\end{lemma}

\begin{proof}
Choose local frames of $B$ and a trivialization of $B^p$.
Let $u_i$ be the resulting units. Their logarithmic differentials
$du_i/u_i$ agree, since the ratios $u_i/u_j$ are $p$th powers.
They define a global K\"ahler form and, in particular, a rational
form $\theta$ that is regular on the smooth locus and closed there.
If this rational form were zero, each $u_i$ would have a $p$th root
in $k(S)$: over a perfect ground field, the kernel of the field
differential is $k(S)^p$. Each root and its inverse are integral over
the normal local rings, hence are regular units. Uniqueness of the
roots makes them compatible with the transition functions of $B$,
trivializing $B$. This is a contradiction.

The restriction of $\theta$ to $\Omega_C^1$ is zero. Indeed, after
trivializing $B|_C$, its $p$th-power trivialization is a global unit
on the proper integral curve, hence a constant. In the conormal
sequence its remaining component would be a section of
$\OO_C(-C)$, which has negative degree. It too is zero.
Thus the exact order $m$ is positive and finite.

Division by the $m$th power of a local equation of $C$ gives a
regular section of $\Omega_S^1(-mC)$ on a neighborhood of $C$.
This assertion holds at every point of $C$: in a regular local
ring, divisibility by a power of the equation of a smooth prime
divisor can be checked at its generic point. The resulting leading
restriction is nonzero. In
\[
0\longrightarrow\OO_C(-(m+1)C)
\longrightarrow\Omega_S^1(-mC)|_C
\longrightarrow\Omega_C^1(-mC)\longrightarrow0,
\]
the first term has negative degree. Therefore the tangential image
is nonzero. Locally write
\[
 \theta=h^m(A\,dh+B_1\,dx),\qquad C=(h=0),
\]
with $B_1$ nonzero at the generic point of $C$.
The coefficient of order $m-1$ in $d\theta$ is
$mB_1h^{m-1}dh\wedge dx$. Closedness gives $p\mid m$.
The nonzero tangential section has a line bundle of degree
$2g(C)-2-mC^2$, proving the inequality. If its degree is zero,
the section has no zeros.
\end{proof}

We will also use the conormal argument without assuming that the
one-form is logarithmic or closed. If a regular one-form has exact
order $j\geq0$ along a smooth proper curve of positive
self-intersection, its nonzero leading restriction has nonzero
tangential image. Otherwise it would give a nonzero section of
$\OO_C(-(j+1)C)$, a line bundle of negative degree.

\begin{lemma}[Resolution loss]\label{lem:positive-square-resolution-loss}
Let $N$ be a normal integral projective surface over an algebraically
closed field, and let $D$ be a reduced integral effective Cartier
divisor. Let $\overline D$ be its normalization, and put
\[
 \Delta_D=p_a(D)-g(\overline D).
\]
There is a projective resolution $f:T\to N$ for which the strict
transform $C_T$ of $D$ is smooth and
\[
 0\leq D^2-C_T^2\leq4\Delta_D.
\]
Here $C_T\simeq\overline D$. The assertion concerns a suitable
resolution, not every resolution obtained by further blowups.
\end{lemma}

\begin{proof}
Choose a projective resolution of $N$, an isomorphism over its
regular locus; such a resolution exists for surfaces of finite type
over a field \cite[Tag~0BGP]{Stacks}. Contract an exceptional smooth
rational curve of self-intersection $-1$ whenever one remains.
The contraction is projective, its target is smooth, and the map to
$N$ factors through it
\cite[Tags~0C5J, 0C2M and~0C2N]{Stacks}.
Each contraction removes an exceptional component, so the process
terminates. Denote the resulting resolution by $f_0:T_0\to N$.

For every exceptional integral curve $E_i$ on $T_0$, the Hodge index
theorem gives $E_i^2<0$: its intersection with the pullback of an
ample divisor on $N$ is zero, whereas that pullback has positive
square. Adjunction gives
\[
 K_{T_0}\cdot E_i=2p_a(E_i)-2-E_i^2\geq0.
\]
Indeed, the only possible negative value would require
$p_a(E_i)=0$ and $E_i^2=-1$. An integral proper curve of arithmetic
genus zero over an algebraically closed field is smooth rational,
and such an exceptional curve was contracted. Consequently
$K_{T_0}\cdot E\geq0$ for every effective exceptional divisor $E$.

Write
\[
 f_0^*D=D_0+E,
\]
where $D_0$ is the strict transform and $E$ is effective exceptional.
Since a pulled-back Cartier divisor has zero intersection with every
exceptional curve, $E^2=-D_0\cdot E$. Therefore
\begin{equation}\label{eq:positive-square-minimal-loss}
 \lambda_0:=D^2-D_0^2=D_0\cdot E\geq0.
\end{equation}

We check that $p_a(f_0^*D)=p_a(D)$, including the contribution of
nonrational surface singularities. Normality gives
$f_{0*}\OO_{T_0}=\OO_N$. The sheaf
$R^1f_{0*}\OO_{T_0}$ is supported on finitely many points, and
$R^if_{0*}\OO_{T_0}=0$ for $i>1$, since the fibres have dimension
at most one. Projection formula gives
\[
 R^1f_{0*}\OO_{T_0}(-f_0^*D)
   \simeq R^1f_{0*}\OO_{T_0}\otimes\OO_N(-D).
\]
Tensoring a finite-support sheaf by an invertible sheaf preserves
its length and hence its Euler characteristic. The two Leray
Euler-characteristic formulas consequently give
\begin{align*}
 \chi(\OO_{f_0^*D})
 &=\chi(\OO_{T_0})-\chi(\OO_{T_0}(-f_0^*D))\\
 &=\chi(\OO_N)-\chi(\OO_N(-D))
 =\chi(\OO_D).
\end{align*}
This proves the asserted arithmetic-genus equality, also when
$f_0^*D$ is nonreduced.

Set $\Delta_0=p_a(D_0)-g(\overline D)\geq0$.
Adjunction on the smooth surface $T_0$, applied to $D_0$ and to
$f_0^*D$, now yields
\begin{equation}\label{eq:positive-square-genus-loss}
 2(\Delta_D-\Delta_0)
    =\lambda_0+K_{T_0}\cdot E\geq\lambda_0.
\end{equation}
Blow up a singular point of the strict transform of $D_0$ whenever
one remains. For one such blowup $\sigma:U'\to U$, let $D_U$ be
the current strict transform, let $\mu\geq2$ be its multiplicity
at the center, let $F$ be the exceptional curve, and let $D_{U'}$
be the new strict transform. Then
\[
 D_{U'}=\sigma^*D_U-\mu F,\qquad
 K_{U'}=\sigma^*K_U+F,\qquad F^2=-1.
\]
Thus $D_{U'}^2=D_U^2-\mu^2$, and adjunction gives
$p_a(D_U)-p_a(D_{U'})=\mu(\mu-1)/2$.
The normalization and its genus are unchanged.
Thus the nonnegative integer $p_a-g(\overline D)$ strictly decreases,
and the process ends with a smooth strict transform $C_T$.
The total arithmetic-genus drop in this second stage is $\Delta_0$.
At each blowup,
\[
 \mu^2\leq2\mu(\mu-1)
       =4\left(\frac{\mu(\mu-1)}2\right),
\]
so the additional self-intersection loss is at most $4\Delta_0$.
Combining this with \eqref{eq:positive-square-genus-loss} gives
\[
 0\leq D^2-C_T^2
 \leq2(\Delta_D-\Delta_0)+4\Delta_0
 \leq4\Delta_D.
\]
The smooth strict transform is the normalization of $D$, as required.
\end{proof}

\begin{lemma}[Genus of the normalized root curve]
\label{lem:positive-square-root-genus}
Let $S$ be a smooth integral projective surface over an algebraically
closed field $k$ of characteristic $p>0$, and let $C\subset S$ be a
smooth integral curve with $d=C^2>0$. Let $B\ne\OO_S$ have order $p$
and be trivial on $C$. Write $m=pa$ for the exact order of its
logarithmic form, and put
\[
 G=2g(C)-2,\qquad \delta=G-md.
\]
Let $\pi:N\to S$ be the normalization of its $\mu_p$-torsor,
and set $D=\pi^*C$ scheme-theoretically. Then $D$ is a reduced
integral effective Cartier divisor,
\[
 D^2=pd,\qquad g(\overline D)=g(C),
\]
and there is an integer $\ell\geq0$ such that
\begin{equation}\label{eq:positive-square-conductor-genus}
 (p-1)\delta=\ell+2\Delta_D,\qquad
 \Delta_D=p_a(D)-g(C).
\end{equation}
In particular $\Delta_D\leq(p-1)\delta/2$.
\end{lemma}

\begin{proof}
Let $\tau:W\to S$ be the original torsor and $n:N\to W$ its
normalization. A local equation of $W$ is $z^p=u$, with $u$ a unit.
The normality argument in
Lemma~\ref{lem:positive-square-logarithmic} shows that $u$ is not a
$p$th power in $k(S)$. Thus $W$ is integral: its locally free
algebra embeds in the degree-$p$ purely inseparable function field.
Normalization is finite for varieties over a field
\cite[Tag~0BXR]{Stacks}, so $N$ is normal, integral and projective.
The finite map $\pi$ is radicial. Indeed, on a normal affine base
the $p$th power of an integral element of $k(N)$ belongs to $k(S)$
and is integral over the base ring, hence belongs to that ring.
This implies uniqueness of the prime above each prime and purely
inseparable residue extensions.

We first compute at the generic point of $C$, keeping its imperfect
residue field. On a sufficiently small smooth neighborhood choose
\'etale coordinates $h,x$, with $C=(h=0)$. If its ring is $R_0$,
the monomials $h^ix^j$, $0\leq i,j<p$, form a basis of $R_0$
over $R_0^p$, by the relative Frobenius isomorphism for \'etale
morphisms \cite[Tag~0EBS]{Stacks}. Write
\[
 u=\sum_{0\leq i,j<p}r_{ij}^{\,p}h^ix^j.
\]
Since $du$ is divisible by $h^{pa}$, comparison of the $dh$
coefficient detects every $i>0$, and comparison of the $dx$
coefficient detects every $j>0$. The factor $h^{pa}$ lies in
$R_0^p$, so freeness gives
$h^{pa}\mid r_{ij}^{\,p}$ for $(i,j)\ne(0,0)$.
Normality along the prime $(h)$ then gives $h^a\mid r_{ij}$.
Consequently
\begin{equation}\label{eq:positive-square-local-root}
 u=f^p+h^{pa}v,\qquad du=h^{pa}dv.
\end{equation}
The function $f$ is a unit along $C$. The tangential leading
coefficient of $du/u$ is $(v_x/u)\,dx$, so
$d\overline v\ne0$ in $\Omega^1_{k(C)/k}$.

Put $R=\OO_{S,\eta_C}$ and $F=k(C)$. The field $F$ is not being
assumed perfect. Since $k$ is perfect and $F/k$ has transcendence
degree one, $[F:F^p]=p$ and the kernel of
$d:F\to\Omega^1_{F/k}$ is $F^p$. Hence $\overline v\notin F^p$.
In the extension field set $t=(z-f)/h^a$. Then
\[
 t^p=v,\qquad z=f+h^at.
\]
The ring
\[
 \mathcal B=R[t]/(t^p-v)
\]
is a DVR with uniformizer $h$: it is finite free of rank $p$ over
$R$, and its quotient by $h$ is the field
$F(\overline v^{1/p})$. It is therefore the integral closure of
$R$ in the extension field. Its residue field is exactly
$F^{1/p}$, since both extensions of $F$ have degree $p$.
In particular $D$ is generically reduced.

The underlying inverse image of $C$ is irreducible because $\pi$
is radicial. A normal surface is Cohen--Macaulay, so its effective
Cartier divisor $D$ satisfies $S_1$. A generically reduced
$S_1$ scheme is reduced. Thus $D$ is reduced and integral.
Its normalization has function field $F^{1/p}$ and is a
Frobenius twist of $C$; base change by the inverse Frobenius
automorphism of the perfect field $k$ preserves its genus.
Thus $g(\overline D)=g(C)$. The finite projection formula gives
$D^2=pC^2$.

We now compute the arithmetic-genus defect. Multiplication in the
graded algebra of $W$, followed by projection to degree zero, is a
perfect pairing. It identifies its finite dualizing module with
the algebra itself, so
\[
 \omega_W\simeq\tau^*\omega_S.
\]
This pairing is not the field trace. Let
\[
 \mathfrak c
   =\mathcal Hom_{\OO_W}(n_*\OO_N,\OO_W)
\]
be the conductor, viewed as an ideal of $\OO_N$. Finite duality
and the preceding identification give
\[
 \omega_N\simeq\pi^*\omega_S\otimes\mathfrak c.
\]
Because $N$ is a normal surface, $\omega_N$ is a rank-one
reflexive sheaf. The same is true of $\mathfrak c$.

At the generic point of $D$, the original torsor order
$\mathcal A=R[z]$ and its normalization have the respective
$R$-bases
\[
 \mathcal A=\bigoplus_{i=0}^{p-1}Rh^{ai}t^i,
 \qquad
 \mathcal B=\bigoplus_{i=0}^{p-1}Rt^i.
\]
The first equality follows by using the basis
$1,z-f,\ldots,(z-f)^{p-1}$ of $R[z]$.
Since $\mathcal B$ is a DVR with uniformizer $h$, its conductor
into $\mathcal A$ is a power of $h$. The displayed bases show that
$h^b\mathcal B\subset\mathcal A$ if and only if
$b\geq a(p-1)$. Thus
\[
 \mathfrak c_{\eta_D}=h^{a(p-1)}\mathcal B.
\]

Define the rank-one reflexive sheaf
\[
 J=\mathfrak c\otimes\OO_N(a(p-1)D).
\]
At the generic point of $D$ it is the unit ideal; at every other
codimension-one point it is an ideal because $\mathfrak c$ is.
Reflexivity on the normal surface therefore gives an inclusion
$J\subset\OO_N$. Put $Q=\OO_N/J$.
At every closed point, both $J$ and $\OO_N$ have depth two.
The depth lemma shows that $Q$ has depth at least one wherever it
is nonzero. Hence $Q$ has no associated point of codimension two.
It has no associated component equal to $D$, since $J$ is the unit
ideal generically on $D$. A local equation of $D$ is consequently a
nonzerodivisor on $Q$. Tensoring
$0\to J\to\OO_N\to Q\to0$ with $\OO_D$ gives an injection
\[
 J|_D\longrightarrow\OO_D.
\]
Its cokernel is supported at finitely many points. Denote its
length by $\ell\geq0$.

Adjunction for a Cartier divisor on a Cohen--Macaulay surface gives
\[
 \omega_D\simeq H\otimes J|_D,\qquad
 H=\pi^*\!\left(\omega_S((1-a(p-1))C)\right)|_D.
\]
Here $H$ is invertible, although $\omega_D$ need not be.
Duality on the integral projective curve $D$ gives
$\chi(\omega_D)=-\chi(\OO_D)$.
The injection just proved, tensored with $H$, has cokernel of
length $\ell$. Riemann--Roch for the invertible sheaf $H$ therefore
yields
\begin{align*}
 2p_a(D)-2
 &=\deg H-\ell\\
 &=p\bigl(K_S\cdot C+(1-a(p-1))d\bigr)-\ell\\
 &=pG-m(p-1)d-\ell.
\end{align*}
Since $2p_a(D)-2=G+2\Delta_D$, rearrangement gives
\eqref{eq:positive-square-conductor-genus}.
No ampleness assumption entered this calculation.
\end{proof}

\begin{lemma}[Root transfer after resolution]
\label{lem:positive-square-resolved-transfer}
Under the hypotheses and notation of
Lemma~\ref{lem:positive-square-root-genus}, suppose
$B\simeq A^{\otimes p}$ for a line bundle $A$ on $S$.
There is a projective resolution $f:T\to N$ on which the strict
transform $C_T$ of $D$ is smooth and
\begin{equation}\label{eq:positive-square-transfer-square}
 g(C_T)=g(C),\qquad
 d_T:=C_T^2\geq pd-2(p-1)\delta.
\end{equation}
If $d_T>0$, the bundle $B_T=(\pi f)^*A$ is nontrivial of order $p$,
is trivial on $C_T$, and its logarithmic form has exact order
$a=m/p$. In particular $p\mid a$, and
\begin{equation}\label{eq:positive-square-excess-recurrence}
 0\leq\delta_T:=G-a d_T
       \leq\bigl(1+2(p-1)a\bigr)\delta.
\end{equation}
\end{lemma}

\begin{proof}
Apply Lemma~\ref{lem:positive-square-resolution-loss} to $D$.
Lemma~\ref{lem:positive-square-root-genus} gives
\[
 d_T\geq D^2-4\Delta_D
       \geq pd-2(p-1)\delta
\]
and identifies the normalization genus with $g(C)$.
This proves \eqref{eq:positive-square-transfer-square}.

Use frames of $A$, with transitions $g_{ij}$, and their $p$th
powers as frames of $B$. The torsor coordinates satisfy
$z_i=g_{ij}^p z_j$. They are regular units on the original torsor
and remain units after normalization and resolution. Thus
\[
 \theta_T=\frac{dz_i}{z_i}
\]
is a global regular closed one-form on $T$. It is the logarithmic
form of $B_T$, whose $p$th power is trivialized by the $z_i$.

At the generic point of $C_T$, the equations in
\eqref{eq:positive-square-local-root} give
$t^p=v$ and $z=f+h^at$. The derivative $v_x$ is a unit there.
The normalization is smooth on this open set, the resolution is an
isomorphism at the generic point of $D$, and $dh,dt$ are a
cotangent basis. The derivation dual to $dt$ kills the original
base ring: it kills $h$, and differentiation of $t^p=v$ then makes
it kill $x$; \'etaleness deals with the remaining functions.
In particular it kills $f$, even when $f$ is not constant.
The coefficient of $dt$ in $\theta_T$ is therefore exactly
\begin{equation}\label{eq:positive-square-transfer-coefficient}
 \frac{h^a}{z}.
\end{equation}
Hence $\theta_T$ is nonzero and its order $j$ along $C_T$ satisfies
$0\leq j\leq a$.

Assume $d_T>0$. If $j<a$, the leading restriction of $\theta_T$
would have zero tangential image: generically the coefficient of
$dt$ has order $a$, while $dh$ is conormal. This would give a
nonzero global section of
$\OO_{C_T}(-(j+1)C_T)$, contradicting its negative degree.
The section is global because divisibility by $h^j$ can be checked
at the generic point in the regular local rings along $C_T$.
Thus $j=a$.

A trivial line bundle on a proper integral variety has zero
logarithmic form for every $p$th-power trivialization, because
global units are constants. Consequently $B_T$ is nontrivial.
Its $p$th power is trivial, so its order is exactly $p$.
Since $a\geq1$, its logarithmic form restricts to zero on $C_T$.
The logarithmic-form map on $p$-torsion line bundles on a smooth
proper curve is injective: a defining unit of zero differential is
a $p$th power in the function field, normality makes the unique
roots local units, and those roots glue. Therefore
$B_T|_{C_T}\simeq\OO_{C_T}$.
Lemma~\ref{lem:positive-square-logarithmic} now gives $p\mid a$
and $\delta_T\geq0$. Finally
\[
 \delta_T
 =G-a d_T
 \leq G-apd+2(p-1)a\delta
 =\bigl(1+2(p-1)a\bigr)\delta,
\]
as asserted.
\end{proof}

The original equality transfer remains a useful special case:
normalization alone already gives the required smooth curve.

\begin{lemma}\label{lem:positive-square-root-transfer}
Under the hypotheses of Lemma~\ref{lem:positive-square-logarithmic},
suppose the tangential leading coefficient is nowhere zero and
$B\simeq A^{\otimes p}$ for an invertible sheaf $A$ on $S$.
Write $m=pq$. Let $\pi:S'\to S$ be the normalization of the
$\mu_p$-torsor associated with $B$.
Then $S'$ is a normal integral projective surface and is smooth
near the smooth integral Cartier divisor $C'=\pi^*C$.
The bundle $B'=\pi^*A$ is nontrivial of order $p$, is trivial on
$C'$, and its logarithmic form $\theta'$ has exact order $q$,
with nowhere-zero tangential leading coefficient. Furthermore,
\[
 (C')^2=pC^2>0,\qquad p\mid q.
\]
\end{lemma}

\begin{proof}
The construction and integrality of the normalized torsor in
Lemma~\ref{lem:positive-square-root-genus} use only normality of
the ambient surface; the local calculation uses only its smooth
neighborhood of $C$. Thus they apply here even when $S$ is singular
away from $C$. At every point of $C$, not just generically, the
nowhere-zero leading coefficient makes $v_x$ a unit in
\eqref{eq:positive-square-local-root}. The finite algebra
\[
 R_0[t]/(t^p-v),\qquad z=f+h^qt,
\]
is then smooth near $h=0$, and its divisor $h=0$ is smooth and
reduced. Its function field is the required extension. After
shrinking around that divisor it is normal, hence is the full
integral closure of $R_0$. These charts identify it with $S'$
near the scheme-theoretic inverse image $C'$. Radiciality makes
$C'$ integral. Its map to $C$ has degree $p$, and its normal
bundle is the pullback of the normal bundle of $C$; hence
$(C')^2=pC^2$.

The proof of Lemma~\ref{lem:positive-square-resolved-transfer}
now applies directly on the smooth neighborhood of $C'$.
The form $dz/z$ has exact order $q$, is nonzero, and gives the
stated nontrivial order-$p$ bundle trivial on $C'$; closedness
gives $p\mid q$. No resolution away from this neighborhood is
needed for these conclusions. Finally
\eqref{eq:positive-square-transfer-coefficient}, with $a=q$,
holds in each of the smooth charts along $C'$. After division by
$h^q$ its tangential coefficient is $1/z$, a unit. It is therefore
nowhere zero, proving the final assertion.
\end{proof}

\begin{proof}[Proof of Theorem~\ref{thm:positive-excess-picard}]
Lemma~\ref{lem:positive-square-logarithmic} gives $m>0$, $p\mid m$
and $\delta=G-md\geq0$, where $G=2g(Y)-2$ and $d=Y^2$.
In particular $r=v_p(m)\geq1$.

First suppose that $M=P^{\otimes p^r}$ for a line bundle
$P$ on $X$. Resolve $X$ by a projective resolution that is an
isomorphism near $Y$. The curve, its self-intersection, its genus,
and the logarithmic order do not change. Pullback on Picard groups
is injective: if $f^*L$ is trivial, normality gives
$f_*\OO=\OO_X$ and projection formula gives $L\simeq\OO_X$.
Thus the pulled-back $M$ is still nontrivial, and the chosen root
$P$ pulls back as well. We may therefore assume $X$ smooth.

Set
\[
 m_i=\frac{m}{p^i},\qquad
 Q_0=1,\qquad
 Q_i=\prod_{j=1}^{i}\bigl(1+2(p-1)m_j\bigr)
 \quad(1\leq i\leq r).
\]
Suppose for contradiction that $Q_r\delta<G$.
We construct successive data $X_i,Y_i,M_i$ with
$X_0=X$, $Y_0=Y$ and $M_0=M$. Write
\[
 d_i=Y_i^2,\qquad \delta_i=G-m_i d_i.
\]
At stage $i$, the bundle $M_i$ is the pullback of
$P^{\otimes p^{r-i}}$ through the composite map to $X$.
For $i<r$ it therefore has the $p$th root given by the pullback
of $P^{\otimes p^{r-i-1}}$.

Inductively assume $d_i>0$, that the logarithmic order is $m_i$,
and that $0\leq\delta_i\leq Q_i\delta$.
Apply Lemma~\ref{lem:positive-square-resolved-transfer} using that
chosen root. Its resolution and smooth strict transform satisfy
\begin{align*}
 m_{i+1}d_{i+1}
 &\geq G-\bigl(1+2(p-1)m_{i+1}\bigr)\delta_i\\
 &\geq G-Q_{i+1}\delta\\
 &\geq G-Q_r\delta>0.
\end{align*}
Thus the next distinguished curve still has positive
self-intersection. The lemma gives a nontrivial order-$p$ bundle
$M_{i+1}$ trivial on it, with exact logarithmic order $m_{i+1}$,
the same genus, and
\[
 0\leq\delta_{i+1}
 \leq\bigl(1+2(p-1)m_{i+1}\bigr)\delta_i
 \leq Q_{i+1}\delta.
\]
This proves the induction. At the final stage,
Lemma~\ref{lem:positive-square-logarithmic} forces $p\mid m_r$.
But $m_r=m/p^r$ is prime to $p$ by the definition of $r$.
The contradiction proves
\[
 Q_r(m)\delta\geq G
\]
whenever $M$ is a $p^r$th tensor power.

It remains to lift divisibility in the N\'eron--Severi group.
For a normal integral projective variety over $k$, the reduced
identity component of the Picard scheme is an abelian variety,
and its $k$-points are the algebraically trivial line-bundle classes
\cite[Theorems~2.5 and~5.4, Proposition~5.10 and
Remark~5.25, pp.~18--20, 37--38, 41--42 and~51]{KleimanPicard}.
Multiplication by $p^r$ is surjective on these points
\cite[Tag~03RP, parts~(3)--(4)]{Stacks}.
If $[M]=p^r[L_0]$ in $\operatorname{NS}(X)$, then
$M\otimes L_0^{-p^r}$ is algebraically trivial. Choose an
algebraically trivial line bundle $L_1$ with
\[
 L_1^{\otimes p^r}\simeq M\otimes L_0^{-p^r}.
\]
Then $M\simeq(L_0\otimes L_1)^{\otimes p^r}$, so the inequality
just proved applies. Since $G=md+\delta$, it is equivalent to
\[
 \bigl(Q_r(m)-1\bigr)\delta\geq md.
\]
An algebraically trivial $M$ has zero N\'eron--Severi class and
satisfies the divisibility hypothesis. This proves the product bound.
The strictness assertion for $r=1$ and odd $p$ is proved in
Theorem~\ref{thm:odd-positive-excess-strict} below, using this weak
bound and its equality cases.
\end{proof}

\begin{corollary}\label{thm:positive-square-picard}
Let $X$ be a normal integral projective surface over an algebraically
closed field $k$ of characteristic $p>0$, and let $Y\subset X$ be a
smooth integral effective Cartier divisor contained in the smooth
locus, with $Y^2>0$. Let $M\ne\OO_X$ have order $p$ and satisfy
$M|_Y\simeq\OO_Y$. If $m$ is the exact order along $Y$ of its
logarithmic form, then
\[
 m\geq1,\qquad p\mid m,\qquad 2g(Y)-2\geq mY^2.
\]
If equality holds and $r=v_p(m)$, then
\[
 [M]\notin p^r\operatorname{NS}(X).
\]
In particular, every algebraically trivial such $M$ satisfies
\[
 2g(Y)-2>mY^2.
\]
\end{corollary}

\begin{proof}
The first assertions are part of
Theorem~\ref{thm:positive-excess-picard}; its quantitative inequality
makes $\delta=0$ impossible when
$[M]\in p^r\operatorname{NS}(X)$, since $mY^2>0$.
\end{proof}

When $v_p(m)=1$, write $m=pb$ with $p\nmid b$.
Then $Q_1(m)-1=2(p-1)b$, and
Theorem~\ref{thm:positive-excess-picard} gives the particularly simple
necessary condition
\begin{equation}\label{eq:positive-square-first-valuation}
 \delta\geq\frac{p}{2(p-1)}Y^2
\end{equation}
for $[M]\in p\operatorname{NS}(X)$, in particular for algebraically
trivial $M$. In characteristic two this says $\delta\geq Y^2$.
The argument uses only the finite root chosen on the original
surface; it requires no Picard divisibility on an intermediate
surface. It imposes a quantitative restriction on positive excess,
but does not decide the existence of a kernel class beyond that
bound.

